\scanpage{40}
\noindent\makebox[0pt][r]{\rule[-1.4ex]{.4pt}{3.7ex}\hspace{.5em}}\uline{Proposition} Soi\deleted{\unsure{ent|t}}t $(C,\sigma,<)$ une \unsure{maquette|moquette} ordonnée. Alors pour tt $x\in C$, les \deleted{\unsure{ens.|sous}} ens.
\[
\begin{aligned}
 C_{\geq x}&=\{y\in C\mid y\geq x\}\\
 C_{\leq x}&=\{y\in C\mid y\leq x\}
\end{aligned}
\]
\makebox[0pt][r]{\rule[-.4ex]{.4pt}{2.4ex}\hspace{.5em}}sont \uline{tot\textsuperscript{\unsure{t|l}}} ordonnés.

Il suffit de \inserted{le} prouver \inserted{pour} \unsure{\deleted{le}|le|\deleted{la}} premier \raisebox{1.3ex}{\unsure{.}\deleted{\unsure{notons|noter}}} (par dualité), \deleted{Soit \unsure{donc}} \unsure{cel|cela} signifie \unsure{aussi|ceci} :
\[
\begin{aligned}
 &\text{Si}\quad x<y,\quad x<z,\quad \text{avec}\quad y\ne z\\
 &\hspace{1em}\text{alors}\quad (y,z)\ \text{comparables.}
\end{aligned}
\]

\noindent\begin{minipage}[c]{.38\linewidth}
\diagram{\begin{tikzcd}[ampersand replacement=\&,row sep={9mm,between origins},column sep={15.3mm,between origins},
execute at end picture={
\draw[line width=.6pt] ([xshift=-16.3mm,yshift=-1.7mm]Rtwofortyx.center) -- ([xshift=-10.9mm,yshift=.4mm]Rtwofortyx.center);
\draw[line width=.45pt] ([xshift=-4mm,yshift=2.4mm]Rtwofortyx.center) -- ([xshift=-2.1mm,yshift=2.9mm]Rtwofortyx.center);
\draw[->,line width=.6pt] ([xshift=5.7mm,yshift=5.9mm]Rtwofortyx.center) -- ([xshift=10.3mm,yshift=8.1mm]Rtwofortyx.center);
\draw[line width=.75pt] ([xshift=7.7mm,yshift=2.1mm]Rtwofortyx.center) -- ([xshift=16.1mm,yshift=-1.3mm]Rtwofortyx.center);
\draw[->,line width=.85pt] ([xshift=1.6mm,yshift=3.7mm]Rtwofortyx.center) -- ([xshift=-10.9mm,yshift=8.3mm]Rtwofortyx.center);
}]
 |[inner sep=0pt,label=above left:{\(z\)}]| \phantom{\cdot} \& \cdot \& |[inner sep=0pt,label=right:{\(y\)}]| \cdot\\
 \& |[alias=Rtwofortyx,inner sep=0pt,label=below:{\(x\)}]| \cdot\arrow[ul,shorten <=4.6mm,shorten >=.8mm]\arrow[ur,shorten <=3.4mm,shorten >=1.9mm] \&
\end{tikzcd}}
\end{minipage}\hfill\begin{minipage}[c]{.57\linewidth}
[On note que $y$ et $z$ sont non congrus mod $\sigma$, à cause de 28 a),]
\end{minipage}

Si $\deleted{\unsure{x}}y$ et $\deleted{\unsure{y}}z$ étaient incomparables, on aurait que $\sigma y$ et $z$ sont comparables, On \unsure{peu|pas|ne} peut avoir $z\unsure{<|\leq}\sigma y$ car on. aurait $x<\sigma y$, ce qui contredit $x<y$ par (28) a). On \deleted{\unsure{ne peut avoir}} \inserted{aurait donc} $\sigma y<z$ d'où $\sigma z<\unsure{\backslash|x}$

% archive sheet number: 39
